\geometry{
  a4paper,
  left=1.8cm,
  right=1.8cm,
  top=2.6cm,
  bottom=2.0cm,
  headheight=26pt,
  headsep=14pt
}

% Set Matmaila (Warm Parchment Clay) Page Background
\pagecolor{bpmatmaila}

% Classic Dictionary Header & Footer
\pagestyle{fancy}
\fancyhf{}
\fancyhead[L]{\small\bfseries\color{bpbluedark}\leftmark}
\fancyhead[R]{\small\bfseries\color{bpbluedark}\rightmark}
\fancyfoot[C]{\small\color{bpgray}\thepage}
\renewcommand{\headrulewidth}{0.5pt}
\renewcommand{\headrule}{\hbox to\headwidth{\color{bpblue}\leaders\hrule height \headrulewidth\hfill}}

% Dictionary Letter Heading Macro (e.g. Large Centered 'अ')
\newcommand{\dictchar}[1]{%
  \vspace{0.8em}%
  \markboth{#1}{#1}%
  {\centerline{\Huge\bfseries\color{bpblue} #1}}%
  \vspace{0.2em}%
  {\color{bpblue}\hrule height 1pt}%
  \vspace{0.8em}%
}

% Star Priority System Macros (TeX Math Stars for Font Compatibility)
\newcommand{\starsThree}{{\small\color{bpblue}$\star\star\star$}}
\newcommand{\starsTwo}{{\small\color{bpgray}$\star\star$}}
\newcommand{\starsOne}{{\small\color{bpbluedark}$\star$}}

\newcommand{\freqcommon}{\starsThree}
\newcommand{\freqformal}{\starsTwo}
\newcommand{\freqrare}{\starsOne}

% Main Dictionary Entry Macro (Kept unbroken on single page/column)
% #1: Devanagari Headword
% #2: Kaithi Script Macros
% #3: Transliteration
% #4: Grammatical Tag & Origin/Dialect
% #5: Star Rating (\starsThree / \starsTwo / \starsOne)
% #6: Hindi Meaning / Context
% #7: Bhojpuri Equivalent / Phrasing
% #8: English Gloss
\newcommand{\bhojentry}[8]{%
  \par\nopagebreak\samepage
  \markright{#1}%
  {\bfseries\large #1} \, #2 \, \textit{\small (#3)} \ \textbf{\small [#4]} \ {#5} \\*
  {\small\textbf{हिन्दी:} #6} \\*
  {\small\textbf{भोजपुरी:} #7} \\*
  {\small\textbf{Eng:} #8}
  \par\vspace{0.5em}%
}

% Hierarchical Sub-Entry Macro (Direct Derivatives with ↳)
\newcommand{\bhojsubentry}[8]{%
  \par\nopagebreak\samepage\hspace*{0.8em}%
  {\bfseries\small $\hookrightarrow$ #1} \, {\small #2} \, \textit{\scriptsize (#3)} \ \textbf{\scriptsize [#4]} \ {#5} \\*
  \hspace*{0.8em}{\small\textbf{हिन्दी:} #6} \\*
  \hspace*{0.8em}{\small\textbf{भोजपुरी:} #7} \\*
  \hspace*{0.8em}{\small\textbf{Eng:} #8}
  \par\vspace{0.4em}%
}

% Compound Word Macro (Double-Indented with →)
\newcommand{\bhojcompound}[7]{%
  \par\nopagebreak\samepage\hspace*{1.6em}%
  {\small\color{bpbluedark}$\rightarrow$ \textbf{#1}} \, {\small #2} \, \textit{\scriptsize (#3)} \ \textbf{\scriptsize [#4]} \\*
  \hspace*{1.6em}{\small\textbf{अर्थ:} #5 / #6} \\*
  \hspace*{1.6em}{\small\textbf{Eng:} #7}
  \par\vspace{0.3em}%
}

% Example Sentence Macro (Mandatory for ★★★)
\newcommand{\bhojexample}[2]{%
  \par\nopagebreak\samepage\vspace{-0.3em}\noindent\hspace*{0.8em}%
  {\small\color{bpbluedark}\textbf{उदा.:} ``#1'' \textit{(#2)}}%
  \par\vspace{0.2em}%
}

% Cross-Reference Macro (देखें: ...)
\newcommand{\bhojsee}[1]{%
  \par\nopagebreak\samepage\vspace{-0.3em}\noindent\hspace*{0.8em}%
  {\small\color{bpgray}\textbf{देखें:} #1}%
  \par\vspace{0.2em}%
}

% Section formatting
\titleformat{\section}
  {\Large\bfseries\color{bpblue}}
  {}{0em}{}
  [\vspace{-0.4em}{\color{bpblue}\rule{\textwidth}{0.8pt}}]

\titleformat{\subsection}
  {\large\bfseries\color{bpbluedark}}
  {\thesubsection}{0.5em}{}

% Paragraph formatting
\setlength{\parskip}{0.3em}
\setlength{\parindent}{0em}
